\section{Post-Freeze Cross-Schema Mechanism Stress Test}
\label{sec:cross_schema}
\begin{table*}[t]
\centering
\caption{Cross-Schema Mechanism Stress Test across Four Heterogeneous Enterprise Domains ($N=32$ Tasks total, 8 per Schema)}
\label{tab:cross_schema}
\setlength{\tabcolsep}{3.5pt}
\resizebox{\textwidth}{!}{%
\begin{tabular}{llcccccccc}
\toprule
\textbf{Enterprise-Style Tool Schema} & \textbf{Domain / Key Constraints} & \textbf{Floor $C_I^*$} & \textbf{Mean $S_0$} & \textbf{Full $S$} & \textbf{Workload Savings} & \textbf{1st-Pass (Exp/Amb)} & \textbf{Final Struct.} & \textbf{Semantic} & \textbf{Ambiguous Recovery} \\
\midrule
\texttt{create\_purchase\_order} & Procurement (Regex, enums, nested items) & 339 & 371 & 552 & \textbf{21.3\%} & 50.0\% (4/4 exp, 0/4 amb) & \textbf{100.0\%} & 87.5\% & \textbf{100.0\%} (4/4) \\
\texttt{inventory\_adjustment} & Warehouse Stock (Signed deltas, reason codes) & 299 & 354 & 541 & \textbf{22.4\%} & 50.0\% (4/4 exp, 0/4 amb) & \textbf{100.0\%} & 87.5\% & \textbf{100.0\%} (4/4) \\
\texttt{shipment\_creation} & Freight Logistics (Nested packages, tracking regex) & 337 & 402 & 646 & \textbf{25.2\%} & 50.0\% (4/4 exp, 0/4 amb) & \textbf{100.0\%} & 100.0\% & \textbf{100.0\%} (4/4) \\
\texttt{customer\_refund} & Financial Accounting (Bounds, currency enums) & 292 & 318 & 474 & \textbf{20.5\%} & 50.0\% (4/4 exp, 0/4 amb) & \textbf{100.0\%} & 100.0\% & \textbf{100.0\%} (4/4) \\
\midrule
\textbf{Unweighted Macro Mean} & Four-Domain Average ($N=4$ domains) & \textbf{316.8} & \textbf{361.3} & \textbf{553.3} & \textbf{22.3\%} & \textbf{50.0\%} (16/16, 0/16) & \textbf{100.0\%} & \textbf{93.8\%} & \textbf{100.0\%} (16/16) \\
\textbf{Pooled Workload Aggregate} & Total Workload Tokens (32 independent tasks, 64 task-arm records) & -- & -- & -- & \textbf{22.5\%} & \textbf{50.0\%} (16/16, 0/16) & \textbf{100.0\%} & \textbf{93.8\%} & \textbf{100.0\%} (16/16) \\
\bottomrule
\end{tabular}%
}
\vspace{1mm}
\begin{minipage}{\textwidth}
\scriptsize
\textit{Note}: Evaluated across 32 independent tasks (4 schemas $\times$ 8 tasks) under a deterministic simulation-based mechanism stress test using whitespace/punctuation heuristic token estimates. Workload savings measure total token reduction relative to monolithic full schema. The \textbf{Unweighted Macro Mean} (22.3\%) averages unrounded per-domain percentage reductions $(21.2691\% + 22.3935\% + 25.1936\% + 20.4859\%) / 4 = 22.3355\%$, while the three new domains average 22.7\% ($(22.3935\% + 25.1936\% + 20.4859\%) / 3 = 22.691\%$). The \textbf{Pooled Workload Aggregate} (22.5\%) computes total token reduction across 32 independent tasks evaluated over two arms (64 task-arm records, reducing pooled tokens from 54,657 to 42,369 tokens across 4 schemas). In this deterministic simulation, both arms share identical modeled structural and recovery outcomes (16/32 first-pass, 32/32 final structural, 30/32 semantic); reported savings reflect lower token exposure under matched simulated completion outcomes. \textbf{Floor $C_I^*$} denotes the schema invariant token floor; \textbf{Mean $S_0$} denotes the actual initial active schema size across evaluated tasks (macro mean: 361.3).
\end{minipage}
\end{table*}


To evaluate cross-schema mechanism transfer beyond purchase orders, we conducted a deterministic simulation-based stress test across three new synthetic schemas plus primary procurement replication (four schemas total) under JSON Schema Draft 2020-12~\cite{jsonschema202012}: (1)~\textbf{\texttt{create\_purchase\_order}} (Procurement: 552 full / 339 $C_I^*$ tok; regexes, enums, budget bounds, items); (2)~\textbf{\texttt{inventory\_adjustment}} (Warehouse: 541 / 299 tok; signed deltas, facilities, reason codes); (3)~\textbf{\texttt{shipment\_creation}} (Logistics: 646 / 337 tok; carriers, tracking, packages, weights); and (4)~\textbf{\texttt{customer\_refund}} (Finance: 474 / 292 tok; refund limits, currencies, return classifications).

\subsection{Evaluation Methodology and Scrutiny Controls}
For each schema, we constructed 8 pre-specified evaluation tasks ($N=32$ total tasks, evaluated across full-schema control and LASR reactive arms yielding $N=64$ task executions), partitioned into two operational regimes: \textbf{Explicit Tasks} (16 tasks, 4 per schema), where intents furnish required enum tokens and parameters; and \textbf{Ambiguous Tasks} (16 tasks, 4 per schema), intentionally constructed as a failure-induction stress test without explicit cues, starting with lean $C_I^*$ to induce initial violations and test closed-loop reactive recomposition.

Methodological controls and scope: (1) this cross-schema suite is a \textbf{deterministic simulation-based mechanism stress test} evaluating error extraction and AST recomposition algorithms rather than live model-provider API billing; (2) token accounting relies on standard whitespace/punctuation heuristic token estimates rather than live provider tokenization; (3) all live Gemini 3.5 Flash-Lite executions and empirical latency measurements are strictly scoped to the primary 32-task benchmark; (4) all 32 tasks, schemas, and invariant floors were pre-specified prior to execution; and (5) the primary frozen benchmark ($N=32$, commit \texttt{9807340}, SHA-256 \texttt{78c2c88e}) remained completely isolated and untouched.

\subsection{Scope of Transfer: Architecture vs. Predictor}
Transfer scope across schemas: \textbf{1) Architecture Transfer}: Closed-loop reactive recovery---error interception, keyword extraction, and in-tree recomposition---operates without domain-specific retuning across the four evaluated heterogeneous JSON Schema structures. \textbf{2) Schema Foundations}: Each schema instantiates its invariant floor ($C_I^* \in \{339, 299, 337, 292\}$ tok, macro mean: $316.8$) and active dictionary $\mathcal{S}_{\text{specs}}$ (yielding task initial sizes $S_0 \in \{371, 354, 402, 318\}$ tok, macro mean: $361.3$), fixed \textit{a priori}. \textbf{3) Predictor Status}: We evaluate recovery transfer without retraining or retuning the primary predictor; explicit tasks test intent-conditioned rules, while ambiguous tasks test zero-shot recovery from $C_I^*$. This evaluates \textbf{architectural and algorithmic mechanism transfer}, rather than universal live-model predictor generalization across arbitrary unseen production APIs.

\subsection{Empirical Robustness Results}
Table~\ref{tab:cross_schema} summarizes system performance across all four enterprise-style domains:

\textit{Domain Token Savings}: Across the three new synthetic domains, LASR consistently reduces cumulative simulated token consumption: saving \textbf{22.4\%} on \texttt{inventory\_adjustment}, \textbf{25.2\%} on \texttt{shipment\_creation}, and \textbf{20.5\%} on \texttt{customer\_refund} (three-domain macro mean: \textbf{22.7\%}, unrounded $22.691\%$), alongside \textbf{21.3\%} on procurement replication. Across all four domains, the unweighted macro mean saving is \textbf{22.3\%} (unrounded $22.3355\%$), and the pooled workload reduction across all 64 per-task records is \textbf{22.5\%} (reducing simulated workload from $54,657$ to $42,369$ tokens). On multi-turn ambiguous tasks requiring reactive disambiguation, a compact foundation ($C_I \approx 290$--$340$ tok) plus targeted reactive delta ($\Delta S \approx 40$ tok) consumes substantially fewer tokens than transmitting monolithic schemas across every turn.

\textit{Matched Recovery Under Lower Token Exposure}: In this deterministic simulation, both full-schema control and reactive LASR exhibit identical modeled structural and recovery dynamics: a 50.0\% first-pass rate (16/16 explicit, 0/16 ambiguous) reflecting the 50/50 stress-test design, followed by a 100.0\% recovery rate on Attempt~2 for the 16 ambiguous tasks, yielding 100.0\% final structural compliance (32/32) and 93.8\% semantic correctness (30/32). Consequently, the empirical finding is not an exclusive recovery capability of LASR over full schema in this simulation, but rather that LASR achieves matched completion fidelity under a 22.5\% reduction in heuristic token exposure.

\textit{High Semantic Fidelity}: Evaluated via independent deterministic domain validators checking operational fidelity (routing, quantities, carrier assignment, currency), LASR attained \textbf{87.5\%} to \textbf{100.0\%} semantic correctness across domains (macro aggregate: \textbf{93.8\%}), verifying that compact schema exposure preserves domain execution fidelity.

These findings provide preliminary simulation evidence of cross-schema transfer: LASR's core architectural mechanics transfer effectively across diverse enterprise-style JSON Schema contracts without retuning the recovery algorithm or the primary predictor; schema-specific constraint vocabularies and invariant foundations were instantiated for each schema.
